\chapter{Implementation, Results and Discussions}

\section{Introduction}
This chapter reports how the four detector configurations of Chapter~III were
trained and evaluated, and what the measurements show. It first fixes the
experimental setup and the evaluation metrics, then describes the C2A dataset and
its split. The results are presented in two parts: a quantitative comparison of
the additive ablation on the C2A test set, and a qualitative analysis of where the
recommended model succeeds and fails, followed by a comparison against published
results on the same dataset and the inference-time and lightweight extensions.
Section~\ref{sec:partb} then reports the extended study carried out after the
ablation: an audit of the benchmark split, a change to the label-assignment step
that improves the smallest targets, a second negative architecture result, and the
adaptation of the model to real aerial footage collected for this thesis. The
chapter closes with a check of the stated objectives and a short cost account.

\section{Experimental Setup}
All models were trained and evaluated on a single workstation with an NVIDIA RTX
4070 Ti SUPER (16~GB, Ada, compute capability 8.9), an Intel Core i7-14700K, and
128~GB of RAM, running Windows~11. The software stack was Python~3.11.9,
PyTorch~2.12.0 (CUDA~12.6, cuDNN~9.10), and Ultralytics~8.4.56, with
pycocotools for COCO-style evaluation. Training used a single GPU with no
distributed data parallelism. Table~\ref{tab:setup} lists the training
configuration; it is identical across the four models so that any measured
difference is due to the architecture, not the optimiser.

\begin{table}[tbp]
  \centering
  \caption{Training configuration, held identical across all four configurations.
  Batch size was set to 16 with an out-of-memory retry ladder; the P2-head models
  settled at 8 because the added stride-4 scale raises activation memory.}
  \label{tab:setup}
  \begin{tabular}{@{}l l@{}}
    \toprule
    Setting & Value \\
    \midrule
    Input resolution & $640 \times 640$ (letterbox) \\
    Optimiser & AdamW, $\text{lr}_0 = 0.001$, $\text{lr}_f = 0.01$, cosine schedule \\
    Epochs (max) & 300, early stop on fitness (patience 50) and F2 (patience 40) \\
    Batch size & 16 (baseline, CBAM); 8 (P2-head models, via OOM ladder) \\
    Mixed precision & enabled (AMP) \\
    Augmentation & YOLO11 defaults; \texttt{close\_mosaic} = 10; no copy--paste/mixup \\
    Seed & 0 (deterministic cuDNN) \\
    \bottomrule
  \end{tabular}
\end{table}

\section{Evaluation Metrics}
Because the task is finding people in disaster imagery, recall is weighted more
heavily than precision: a missed victim is costlier than a false alarm that an
operator can dismiss. The primary operating metric is therefore the $F_2$ score,
which weights recall four times as heavily as precision. With true positives
$\mathit{TP}$, false positives $\mathit{FP}$, and false negatives $\mathit{FN}$
(a detection counts as $\mathit{TP}$ when its intersection-over-union with a
ground-truth box is at least $0.5$), precision $P$, recall $R$, and the
$F_\beta$ family are
\begin{align}
  P = \frac{\mathit{TP}}{\mathit{TP}+\mathit{FP}}, \qquad
  R = \frac{\mathit{TP}}{\mathit{TP}+\mathit{FN}}, \qquad
  F_\beta = (1+\beta^{2})\,\frac{P \cdot R}{\beta^{2} P + R},
\end{align}
with $\beta = 1$ for $F_1$ and $\beta = 2$ for $F_2$. Localisation quality is
measured by COCO~\cite{lin2014coco} average precision: $\mathrm{AP}$ averaged over
IoU thresholds $0.5$ to $0.95$, along with $\mathrm{AP}_{50}$, $\mathrm{AP}_{75}$,
and the small-object $\mathrm{AP}_S$. Because the targets are extremely small, recall is
also reported per object-size band (very-tiny $<8$~px, tiny $8$--$16$~px,
small $16$--$32$~px, and medium $32$--$96$~px), which is where the effect of
the P2 head is visible. Efficiency is reported as parameter count, GFLOPs, and
end-to-end latency measured on the 4070 Ti SUPER. AP-style metrics use
$\text{conf} = 0.001$, $\text{IoU} = 0.7$; operational $F_1$/$F_2$ use the
confidence threshold that maximises each score.

\section{Dataset}
The models are trained and tested on C2A (Combination to Application), a dataset
for human detection in disaster scenes~\cite{nihal2024c2a}. C2A is semi-synthetic:
disaster backgrounds from AIDER are combined with human poses segmented from
pose-estimation datasets by a U$^2$-Net model, which are then randomly scaled and
placed and annotated with horizontal bounding boxes. The dataset contains 10{,}215
images and roughly 360{,}000 labelled human instances across four scene types
(collapsed buildings, fire, flood, and traffic accident). Targets are dominated by
very small objects: about 47\% of people occupy fewer than 10~px and about 52\%
fall between 10 and 50~px. All pose sub-labels are collapsed to a single
\texttt{person} class; the C2A annotation stores this class under a numeric
placeholder, which is mapped to \texttt{person} at load time. The frozen split
used throughout this thesis is given in Table~\ref{tab:split}; the split hashes are
recorded so that every model sees exactly the same images.

\begin{table}[tbp]
  \centering
  \caption{C2A split used for all experiments. The same frozen split (verified by
  MD5) is used by every configuration.}
  \label{tab:split}
  \begin{tabular}{@{}l r r r@{}}
    \toprule
    & Train & Validation & Test \\
    \midrule
    Images & 6{,}129 & 2{,}043 & 2{,}043 \\
    \bottomrule
  \end{tabular}
\end{table}

Representative samples covering all four scene types and the range of target
scales are shown in Figure~\ref{fig:dataset_samples}.

\begin{figure}[tbp]
  \centering
  \begin{minipage}{0.44\linewidth}\centering
    \includegraphics[width=\linewidth,height=3.7cm,keepaspectratio]{figures/c2a_collapsed.png}\\[2pt]
    {\small (a) collapsed building}
  \end{minipage}\hfill
  \begin{minipage}{0.44\linewidth}\centering
    \includegraphics[width=\linewidth,height=3.7cm,keepaspectratio]{figures/c2a_fire.png}\\[2pt]
    {\small (b) fire}
  \end{minipage}\\[8pt]
  \begin{minipage}{0.44\linewidth}\centering
    \includegraphics[width=\linewidth,height=3.7cm,keepaspectratio]{figures/c2a_flood.png}\\[2pt]
    {\small (c) flood}
  \end{minipage}\hfill
  \begin{minipage}{0.44\linewidth}\centering
    \includegraphics[width=\linewidth,height=3.7cm,keepaspectratio]{figures/c2a_traffic.png}\\[2pt]
    {\small (d) traffic incident}
  \end{minipage}
  \caption{Representative C2A test samples, one per scene type: (a)~collapsed
  building, (b)~fire, (c)~flood, and (d)~traffic incident. Human targets are small
  and numerous in every scene, which is the central difficulty of the dataset.}
  \label{fig:dataset_samples}
\end{figure}

\section{Implementation and Results}

\subsection{Detector Family Benchmark}
Before the ablation, the YOLOv9~\cite{wang2024yolov9}, YOLOv10~\cite{wang2024yolov10},
and YOLO11~\cite{jocher2024yolo11} families were benchmarked on C2A to select a
baseline; the comparison is given in Table~\ref{tab:detbench}. This benchmark was run
under an earlier protocol, so its absolute values are not directly comparable to the
final ablation of Table~\ref{tab:main}; it is used only to choose the backbone. The
largest model, YOLOv9-e, gives the highest accuracy (mAP50 $0.855$, mAP50-95 $0.619$)
but at $111.8$~MB, about three times the size of the mid-size models. YOLO11m
reaches mAP50 $0.841$ and mAP50-95 $0.592$ at $38.6$~MB, close to the best while
staying small enough to deploy, and gives the best accuracy per megabyte of the
family. It was therefore chosen as the baseline.

\begin{table}[tbp]
  \centering
  \caption{Preliminary detector-family benchmark on C2A (test set) used to select
  the baseline. YOLO11m was chosen for its accuracy-to-size balance.}
  \label{tab:detbench}
  \setlength{\tabcolsep}{5pt}
  \begin{tabular}{@{}l r r r r r r@{}}
    \toprule
    Model & mAP50 & mAP50-95 & $P$ & $R$ & $F_1$ & Size (MB) \\
    \midrule
    YOLOv9-s   & 0.787 & 0.504 & 0.841 & 0.736 & 0.785 & 14.5 \\
    YOLOv9-m   & 0.823 & 0.565 & 0.863 & 0.772 & 0.815 & 38.9 \\
    YOLOv9-e   & \textbf{0.855} & \textbf{0.619} & \textbf{0.884} & \textbf{0.805} & \textbf{0.843} & 111.8 \\
    YOLOv10-s  & 0.792 & 0.517 & 0.836 & 0.738 & 0.784 & 15.8 \\
    YOLOv10-m  & 0.827 & 0.571 & 0.857 & 0.777 & 0.815 & 31.9 \\
    YOLOv10-l  & 0.843 & 0.598 & 0.870 & 0.792 & 0.829 & 49.8 \\
    YOLOv11-s  & 0.803 & 0.535 & 0.845 & 0.757 & 0.798 & 18.3 \\
    \textbf{YOLO11m (selected)} & 0.841 & 0.592 & 0.867 & 0.794 & 0.829 & 38.6 \\
    YOLOv11-l  & 0.847 & 0.603 & 0.876 & 0.796 & 0.834 & 48.8 \\
    \bottomrule
  \end{tabular}
\end{table}

\subsection{Attention Module Selection}
The choice of CBAM over the lighter ECA module~\cite{wang2020eca} was settled by a
preliminary comparison on the baseline, summarised in Table~\ref{tab:eca} (again an
earlier protocol; values are relative, not comparable to Table~\ref{tab:main}).
Replacing the baseline attention with CBAM gave the highest $F_2$ ($0.849$) and the
highest small-object recall ($0.892$, $+1.3$ points over the baseline), whereas ECA
barely moved either metric. CBAM also reduced parameters relative to the baseline
block. CBAM was therefore adopted as the attention module and carried into the main
ablation.

\begin{table}[tbp]
  \centering
  \caption{Preliminary attention-module comparison on C2A (earlier protocol).
  CBAM gives the best $F_2$ and small-object recall, so it is used in the ablation.}
  \label{tab:eca}
  \begin{tabular}{@{}l r r r@{}}
    \toprule
    Model & $F_2$ & Small-object recall & $F_1$ \\
    \midrule
    YOLO11m baseline & 0.844 & 0.881 & 0.837 \\
    + ECA            & 0.844 & 0.884 & 0.840 \\
    + CBAM           & \textbf{0.849} & \textbf{0.892} & \textbf{0.842} \\
    \bottomrule
  \end{tabular}
\end{table}

\subsection{Quantitative Results}
Table~\ref{tab:main} reports the four configurations on the C2A test set under the
COCO protocol. Aggregate average precision is saturated: all four models sit within
$0.002$ AP of one another. The differences appear in the metrics that matter for
this task. Adding the P2 head raises $\mathrm{AP}_{50}$ from $0.843$ to $0.853$,
$\mathrm{AR}_{100}$ from $0.691$ to $0.703$, and the operational $F_2$ from
$0.840$ to $0.844$, at a cost of about $0.5$~M parameters and $1$~ms of latency.
CBAM on its own is close to free: it removes about $1$~M parameters relative to the
baseline (it replaces the two heavier C2PSA blocks) while matching or slightly
improving $F_1$/$F_2$ and giving the lowest latency of the four. The state-space
neck is the outlier in the wrong direction: it adds $2.4$~M parameters and roughly
triples end-to-end latency ($14.6 \to 41.1$~ms) without improving any accuracy
metric. On this evidence the recommended configuration is CBAM~+~P2, and the
state-space neck is not carried forward.

\begin{table}[tbp]
  \centering
  \caption{Additive ablation on the C2A test set (COCO protocol). AP is averaged
  over IoU 0.5--0.95; $F_1$ and $F_2$ are at their optimal confidence thresholds;
  latency is end-to-end on the RTX 4070 Ti SUPER. Best value in each column in bold.}
  \label{tab:main}
  \setlength{\tabcolsep}{4pt}
  \begin{tabular}{@{}l r r r r r r r r r@{}}
    \toprule
    Model & Params & GFLOPs & AP & $\mathrm{AP}_{50}$ & $\mathrm{AP}_{75}$ & $\mathrm{AR}_{100}$ & $F_1$ & $F_2$ & Lat.\ (ms) \\
    & (M) & & & & & & & & \\
    \midrule
    YOLO11m baseline    & 20.03 & 67.7 & \textbf{0.615} & 0.843 & 0.655 & 0.691 & \textbf{0.850} & 0.840 & 13.7 \\
    + CBAM              & 19.08 & 66.9 & \textbf{0.616} & 0.847 & 0.656 & 0.692 & \textbf{0.850} & 0.841 & \textbf{13.5} \\
    + CBAM + P2         & 19.57 & 86.7 & 0.615 & \textbf{0.853} & 0.660 & 0.703 & 0.848 & \textbf{0.844} & 14.6 \\
    + Mamba + CBAM + P2 & 22.01 & 98.4 & 0.614 & 0.852 & \textbf{0.662} & \textbf{0.704} & 0.846 & \textbf{0.844} & 41.1 \\
    \bottomrule
  \end{tabular}
\end{table}

The per-size recall in Table~\ref{tab:persize} explains the aggregate flatness. The
P2 head raises very-tiny ($<8$~px) recall from $0.743$ to $0.757$ ($+1.5$~points),
the band that matters most for spotting distant victims, but it slightly lowers
recall in the tiny, small, and medium bands. The medium band drops most in relative
terms, but it contains only 317 of the roughly 72{,}500 test instances, so its
contribution to the aggregate is small and its estimate is noisy. The net effect is
a model that is better at the smallest targets and roughly even elsewhere, which is
the intended trade for this application. The state-space neck tracks CBAM~+~P2
almost exactly here, confirming it adds nothing in the size regime it was meant to
help.

\begin{table}[tbp]
  \centering
  \caption{Per-size recall on the C2A test set, by target pixel area. Ground-truth
  counts per band are shown; the medium band is small (317 instances) and its
  recall estimate is correspondingly noisy.}
  \label{tab:persize}
  \begin{tabular}{@{}l r r r r@{}}
    \toprule
    Model & very-tiny & tiny & small & medium \\
    & ($<8$px) & ($8$--$16$) & ($16$--$32$) & ($32$--$96$) \\
    GT instances & 25{,}072 & 20{,}520 & 26{,}614 & 317 \\
    \midrule
    YOLO11m baseline    & 0.743 & 0.869 & 0.894 & 0.909 \\
    + CBAM              & 0.746 & 0.873 & 0.895 & 0.912 \\
    + CBAM + P2         & \textbf{0.757} & 0.865 & 0.886 & 0.808 \\
    + Mamba + CBAM + P2 & \textbf{0.757} & 0.870 & 0.887 & 0.811 \\
    \bottomrule
  \end{tabular}
\end{table}

Figure~\ref{fig:curves} shows the precision--recall and $F_1$--confidence curves for
the recommended CBAM~+~P2 model, and Figure~\ref{fig:calib} its reliability
diagram; the expected calibration error is $0.021$, so the confidence scores are
close to their empirical accuracy. The per-epoch training curves are in
Figure~\ref{fig:training}.

\begin{figure}[tbp]
  \centering
  \begin{minipage}{0.49\linewidth}\centering
    \figorplaceholder{figures/cbam_p2_pr_curve.png}{6.0cm}
  \end{minipage}\hfill
  \begin{minipage}{0.49\linewidth}\centering
    \figorplaceholder{figures/cbam_p2_f1_conf.png}{6.0cm}
  \end{minipage}
  \caption{Recommended CBAM~+~P2 model on the C2A test set: (a)~precision--recall
  curve and (b)~$F_1$ score against confidence threshold, from which the operating
  point is chosen.}
  \label{fig:curves}
\end{figure}

\begin{figure}[tbp]
  \centering
  \begin{minipage}{0.49\linewidth}\centering
    \figorplaceholder{figures/cbam_p2_calibration.png}{6.0cm}
  \end{minipage}\hfill
  \begin{minipage}{0.49\linewidth}\centering
    \figorplaceholder{figures/cbam_p2_confusion.png}{6.0cm}
  \end{minipage}
  \caption{CBAM~+~P2 model: (a)~reliability diagram (expected calibration error
  $0.021$) and (b)~confusion matrix on the C2A test set.}
  \label{fig:calib}
\end{figure}

\begin{figure}[tbp]
  \centering
  \figorplaceholder{figures/cbam_p2_training_panel.png}{7.5cm}
  \caption{Per-epoch training and validation dynamics for the CBAM~+~P2 model: box,
  classification, and distribution-focal losses together with precision, recall,
  and mAP over the course of training.}
  \label{fig:training}
\end{figure}

\subsection{Qualitative Results and Error Analysis}
The per-size results in Table~\ref{tab:persize} show where the recommended model
still fails, and inspection of the missed detections groups them into three
recurring types. The first is heavy occlusion, where a person is mostly hidden by
debris or water and only a small fragment stays visible. The second is the
extreme-scale case, where a target spans only a few pixels and carries too little
texture to separate it from the background. The third is crowding, where several
people stand close together and the detector merges them into one box. These types
match the numbers. The very-tiny band holds the largest share of the residual
error, while the medium band, which contains few instances and larger targets, is
almost saturated. The trade the model makes is therefore the intended one for this
task: it recovers more of the smallest targets, where a missed detection is a missed
person, at a small cost on the easier size bands.

\subsection{Comparison with State-of-the-Art}
Table~\ref{tab:sota} places the recommended model against the published results of
Nihal et al.~\cite{nihal2024c2a} on the same C2A test split, none of which use
test-time slicing or augmentation. The CBAM~+~P2 model matches or exceeds every
convolutional and transformer baseline in that study except YOLOv9-e, and it does
so with about $19.6$~M parameters against YOLOv9-e's $57.3$~M, roughly a third of the
size. This compactness is the basis for the deployment discussion in
Section~\ref{sec:lightweight}.

\begin{table}[tbp]
  \centering
  \caption{Comparison with published detectors on the C2A test set. Baseline
  numbers are from Nihal et al.~\cite{nihal2024c2a}, which use standard training
  without slicing or test-time augmentation.}
  \label{tab:sota}
  \begin{tabular}{@{}l r r r@{}}
    \toprule
    Model & $\mathrm{AP}_{50}$ & AP & Params \\
    \midrule
    Faster R-CNN            & 0.634 & 0.366 & $\sim$41M \\
    RetinaNet               & 0.693 & 0.383 & $\sim$37M \\
    RTMDet                  & 0.708 & 0.442 & --- \\
    Cascade R-CNN           & 0.735 & 0.486 & $\sim$69M \\
    DINO                    & 0.789 & 0.471 & $\sim$47M \\
    YOLOv5                  & 0.808 & 0.492 & --- \\
    YOLOv9-e                & \textbf{0.893} & \textbf{0.688} & 57.3M \\
    \midrule
    YOLO11m baseline (ours) & 0.843 & 0.615 & 20.1M \\
    CBAM + P2 (ours)        & 0.853 & 0.615 & \textbf{19.6M} \\
    \bottomrule
  \end{tabular}
\end{table}

\subsection{Inference-Time Enhancement: SAHI and TTA}
Slicing-aided hyper inference (SAHI)~\cite{akyon2022sahi} runs the trained detector
over overlapping crops so that very small targets occupy a larger fraction of each
input tile; test-time augmentation (TTA) instead runs the detector at several
resolutions and merges the results. Both are applied as separate, clearly-labelled
inference-time settings, not folded into the primary table, because the sliced
protocol uses per-image box matching for which the standard mAP is not directly
comparable. Table~\ref{tab:sahi} reports the SAHI sweep. A $256$-px slice gives the
highest very-tiny recall, $0.829$ (up from $0.767$, $+6.2$ points), but at a heavy
cost of $534$~ms per image; a $512$-px slice is the best balance, reaching $F_1$
$0.866$ and very-tiny recall $0.808$ at $187$~ms. TTA at $1280$~px raised mAP50 by
$1.6$ points and mAP50-95 by $5.4$ points with no retraining and cut the person
miss-rate from $14\%$ to $10\%$; beyond $1280$~px it degraded, as a model trained at
$640$~px cannot extrapolate to three times the resolution. These inference-time
results were measured on the state-space variant. Because that variant tracks
CBAM~+~P2 within $0.001$ AP on every accuracy metric (Table~\ref{tab:main}), the
slice-size ranking is expected to carry over; a re-measurement on CBAM~+~P2 was
set aside in favour of the extended study of Section~\ref{sec:partb}, which
addresses the same very-tiny-recall goal at training time instead of at
inference time.

\begin{table}[tbp]
  \centering
  \caption{SAHI slice-size sweep on the C2A test set (per-image box matching at
  IoU~0.5), reported as a separate inference-time setting. Baseline row is the
  same model without slicing.}
  \label{tab:sahi}
  \begin{tabular}{@{}l r r r r r r@{}}
    \toprule
    Setting & Slice & $P$ & $R$ & $F_1$ & $F_2$ & very-tiny recall \\
    \midrule
    no SAHI  & ---  & 0.855 & 0.845 & 0.850 & 0.847 & 0.767 \\
    + SAHI   & 512  & 0.867 & 0.865 & \textbf{0.866} & 0.865 & 0.808 \\
    + SAHI   & 256  & 0.790 & 0.872 & 0.829 & 0.854 & \textbf{0.829} \\
    \bottomrule
  \end{tabular}
\end{table}

\subsection{Lightweight Variant}
\label{sec:lightweight}
The recommended CBAM~+~P2 model is already compact. At $19.6$~M parameters it is
smaller than the plain YOLO11m baseline and about a third the size of the largest
published detector on C2A, and it runs at $14.6$~ms per image on a single consumer
GPU, which is real-time for drone review at typical frame rates. For on-board use on
constrained hardware, the same recipe, meaning the CBAM substitution and the added
P2 scale, transfers directly to a smaller YOLO11 backbone such as YOLO11s, which
trades a little accuracy for a further cut in parameters and latency. Building that
variant and measuring it on an embedded device is set out as future work in
Chapter~VII; the design presented here is already small enough to deploy without it.

\section{Extended Study: Protocol Hardening and Real-World Adaptation}
\label{sec:partb}
The ablation above answers which architectural additions help on the benchmark as
published. This section reports a second phase of work that hardens that answer
and carries the model to real imagery. It has four parts. The benchmark split is
audited for scene leakage and the headline model is re-measured on a
scene-disjoint re-split. Under that protocol, the search for further gains moves
from the network to the label-assignment step, where a similarity that stays
informative at tiny scales raises recall of sub-8-pixel targets; a second
architectural branch, tested the
same way, does not survive its ablation and is reported as a negative result.
Real aerial footage collected and annotated for this thesis then measures how far
the synthetic-trained detector transfers, and what a small budget of real labels
buys, including a held-out test on a disaster event the model has never seen.
Figure~\ref{fig:partbchain} summarises the sequence. All experiments in this
section start from the recommended CBAM~+~P2 configuration and run on the
workstation of Section~4.2; training settings that differ from
Table~\ref{tab:setup} are stated where they apply, and every number in this
section is reproducible from archived run configurations.

\begin{figure}[tbp]
  \centering
  \figorplaceholder{figures/fig_partB_chain.png}{4.5cm}
  \caption{The extended study as a chain of decisions. Starting from the
  recommended CBAM~+~P2 model, the evaluation protocol is hardened with a
  scene-disjoint re-split; a frequency-gated architecture branch fails its
  ablation and is dropped; a tiny-aware change to label assignment is kept; and
  the resulting model is adapted to real imagery with a small labelled budget.}
  \label{fig:partbchain}
\end{figure}

\subsection{Scene Leakage and the Scene-Disjoint Re-Split}
C2A composites several synthetic images from each disaster background, and the
official split assigns those images to training and test sets independently, so
the same background scene can appear on both sides. An audit of scene identifiers
found that 98\% of test-set scenes also occur in the training set. To measure the
effect, the dataset was re-partitioned so that all images from one scene fall on
one side of the split (6{,}135 train, 2{,}040 validation, and 2{,}040 test
images), and the CBAM~+~P2 model was retrained from the same initialisation and
schedule. Table~\ref{tab:leak} compares the two protocols. The clean split costs
$1.6$ points of $\mathrm{AP}_{50}$ and about one point of very-tiny recall, so
the leakage inflates the official numbers only modestly, and the conclusions of
the ablation stand. All experiments in the remainder of this section use the
scene-disjoint split, whose test set holds 72{,}847 person instances (24{,}806
below 8~px, 20{,}893 at 8--16~px, 26{,}738 at 16--32~px, and 410 at 32--96~px).

\begin{table}[tbp]
  \centering
  \caption{Effect of scene leakage: the recommended CBAM~+~P2 model evaluated
  under the official split and retrained under the scene-disjoint re-split. The
  gap measures how much the shared scenes inflate the official protocol.}
  \label{tab:leak}
  \begin{tabular}{@{}l r r r@{}}
    \toprule
    Protocol & $\mathrm{AP}_{50}$ & $\mathrm{AP}_S$ & Very-tiny recall \\
    \midrule
    Official split          & 0.853 & 0.616 & 0.757 \\
    Scene-disjoint re-split & 0.837 & 0.613 & 0.745 \\
    Difference (points)     & $-1.6$ & $-0.2$ & $-1.2$ \\
    \bottomrule
  \end{tabular}
\end{table}

\subsection{Tiny-Aware Label Assignment}
\label{sec:assignresults}
The remaining error is concentrated below 8~px, where the intersection-over-union
between a candidate box and a ground-truth box a few pixels wide is numerically
unstable: a one-pixel shift can drive it to zero, so the assigner that selects
positive training samples ranks candidates for the smallest targets almost at
random. Section~\ref{sec:assignmethod} describes the remedy studied here, which blends a
normalised Wasserstein similarity into the localisation term of the task-aligned
assigner with a weight $\alpha$, leaving the architecture and the loss unchanged.

Fifty-epoch pilot runs swept $\alpha$ over $\{0, 0.5, 1.0\}$; Table~\ref{tab:alpha}
and Figure~\ref{fig:alphadose} show the outcome. The response is monotone in both
directions: raising $\alpha$ buys very-tiny recall ($0.693 \to 0.712 \to 0.729$)
and sells accuracy on larger targets, slowly at first and then steeply
($\mathrm{AP}_S$ $0.560 \to 0.553 \to 0.505$). The full replacement at
$\alpha=1.0$ pays about three points of $\mathrm{AP}_S$ for its recall gain,
which is a poor trade; $\alpha=0.5$ keeps most of the gain at a tenth of the
cost and was carried forward.

\begin{table}[tbp]
  \centering
  \caption{Dose--response of the assignment blend weight $\alpha$ in fifty-epoch
  pilot runs on the scene-disjoint test set. Recall per size band is measured at
  each model's own optimal-$F_1$ threshold.}
  \label{tab:alpha}
  \begin{tabular}{@{}l r r r r r@{}}
    \toprule
    Blend & Very-tiny & Tiny & Small & $\mathrm{AP}_S$ & $\mathrm{AP}_{50}$ \\
    \midrule
    $\alpha=0$ (CIoU only) & 0.693 & 0.812 & 0.830 & 0.560 & 0.821 \\
    $\alpha=0.5$           & 0.712 & 0.809 & 0.820 & 0.553 & 0.823 \\
    $\alpha=1.0$ (NWD only) & 0.729 & 0.795 & 0.797 & 0.505 & 0.808 \\
    \bottomrule
  \end{tabular}
\end{table}

\begin{figure}[tbp]
  \centering
  \figorplaceholder{figures/fig_alpha_dose.png}{7.0cm}
  \caption{The assignment blend weight $\alpha$ acts as a trade dial: very-tiny
  ($<8$~px) recall rises monotonically with $\alpha$ while $\mathrm{AP}_S$ falls,
  slowly to $\alpha=0.5$ and steeply beyond it. The dashed line marks the value
  carried into the full-protocol comparison.}
  \label{fig:alphadose}
\end{figure}

The decisive test is a matched pair at the full 300-epoch protocol: two models
trained by the same launcher from the same pretrained initialisation on the same
split, differing only in $\alpha$ ($0$ against $0.5$). Table~\ref{tab:d2pair}
reports the pair together with confidence intervals from the bootstrap procedure
of Section~\ref{sec:stats}. The assignment change raises very-tiny recall from
$0.730$ to $0.751$ ($+2.08$ points, 95\% CI $[1.77, 2.38]$, $p<0.001$) and
$\mathrm{AP}_S$ by two points, holds the overall operating score ($F_2$ $+0.27$
points, $p=0.024$) and $\mathrm{AP}_{50}$ ($-0.08$ points, not significant), and
pays for the gain with significant losses of about one point in the 8--16 and
16--32~px bands and a larger relative loss in the sparse 32--96~px band. The
mechanism therefore reallocates the detector's capacity toward the smallest
targets rather than improving everything at once. For search and rescue, where
the sub-8-pixel band holds the most easily missed people, this is the intended
direction of trade, and it is the only intervention in this thesis that moved
very-tiny recall by more than run-to-run noise. The matched control reaches an
absolute level about $0.8$ points of $\mathrm{AP}_{50}$ above the re-baseline of
Table~\ref{tab:leak}, which reflects launcher and early-stopping differences
between the two training scripts; the pair in Table~\ref{tab:d2pair} is
internally matched, so its deltas are unaffected.

\begin{table}[tbp]
  \centering
  \caption{Matched full-protocol pair on the scene-disjoint test set: identical
  training except for the assignment blend ($\alpha=0$ against $\alpha=0.5$).
  Deltas carry 95\% bootstrap confidence intervals over test images
  (Section~\ref{sec:stats}); $\mathrm{AP}_S$ is computed by the COCO toolkit and
  is not bootstrapped. The 32--96~px band holds only 410 instances and its
  estimate is correspondingly noisy.}
  \label{tab:d2pair}
  \setlength{\tabcolsep}{4.5pt}
  \begin{tabular}{@{}l r r r c r@{}}
    \toprule
    Metric & $\alpha=0$ & $\alpha=0.5$ & $\Delta$ (points) & 95\% CI & $p$ \\
    \midrule
    $\mathrm{AP}_{50}$    & 0.845 & 0.844 & $-0.08$ & $[-0.32, 0.17]$ & 0.54 \\
    $F_2$ (operational)   & 0.826 & 0.828 & $+0.27$ & $[0.03, 0.50]$  & 0.024 \\
    $\mathrm{AP}_S$       & 0.595 & 0.615 & $+2.00$ & --- & --- \\
    Very-tiny recall      & 0.730 & 0.751 & $+2.08$ & $[1.77, 2.38]$  & $<0.001$ \\
    Tiny recall           & 0.834 & 0.823 & $-1.15$ & $[-1.48, -0.84]$ & $<0.001$ \\
    Small recall          & 0.850 & 0.841 & $-0.95$ & $[-1.40, -0.50]$ & $<0.001$ \\
    Medium recall         & 0.363 & 0.232 & $-13.2$ & $[-17.9, -9.3]$ & $<0.001$ \\
    \bottomrule
  \end{tabular}
\end{table}

\subsection{A Frequency-Gated Branch That Did Not Help}
The assignment result was reached only after a second architecture branch failed.
Motivated by a degradation probe in which near-lossless JPEG re-compression
(quality 90) collapsed $\mathrm{AP}_{50}$ by 65 points, which shows the
tiny-object signal lives in the highest spatial-frequency band, a branch was
built that extracts high-frequency evidence and gates it with learned
large-kernel context at the two finest fusion points of the neck (Section~\ref{sec:fccgmethod}).
The module trains stably and its gates stay active, but in a fifty-epoch paired
ablation against the same control it changed nothing that matters: very-tiny
recall $0.691$ against the control's $0.693$, $\mathrm{AP}_S$ $0.554$ against
$0.560$, at a cost of $0.55$~M parameters and 12 additional GFLOPs. The result
was reproduced on a second machine. As with the state-space neck in the main
ablation, the branch is excluded, and the two negatives together point the same
way: for this task the productive lever is training supervision, not added
architecture.

\subsection{Adaptation to Real Imagery}
\label{sec:realresults}
C2A is composited, so the question that matters for deployment is what happens on
real footage. Three real sets were collected and annotated for this thesis
(Section~\ref{sec:adaptmethod}): drone video recorded at three altitudes over open ground, campus
scenes with partial occlusion, and disaster-scene frames taken from published
news footage of two events, one flood and one building collapse. A fourth set,
frames from a news clip of a third, unseen disaster event, is held out entirely
and appears only in the next subsection. Table~\ref{tab:owndata} lists all sets;
Figure~\ref{fig:realsamples} shows samples. Test frames are frozen and never
trained on, and the disaster training frames were checked for near-duplicate
overlap with every evaluation frame using a perceptual hash before annotation.

\begin{table}[tbp]
  \centering
  \caption{Real-imagery sets collected and annotated for this thesis. Median box
  size is the square root of box area over all person annotations in the set;
  train and test frames never share source footage except in the within-video
  disaster case discussed in the text.}
  \label{tab:owndata}
  \setlength{\tabcolsep}{5pt}
  \begin{tabular}{@{}l l r r r@{}}
    \toprule
    Set & Role & Frames & Persons & Median (px) \\
    \midrule
    Drone, 3 altitudes  & fine-tune train      & 101 & 6{,}313 & --- \\
    Drone, 3 altitudes  & fine-tune validation & 14  & 973     & --- \\
    Drone test          & evaluation (frozen)  & 60  & 3{,}584 & 44 \\
    Campus (occlusion)  & fine-tune train      & 87  & 4{,}510 & 37 \\
    Campus test         & evaluation (frozen)  & 25  & 663     & 37 \\
    Disaster (2 events) & fine-tune train      & 60  & 1{,}341 & 23 \\
    Disaster R-set      & evaluation (frozen)  & 25  & 1{,}070 & 17 \\
    Unseen-event set    & evaluation (frozen)  & 23  & 318     & 28 \\
    \bottomrule
  \end{tabular}
\end{table}

\begin{figure}[tbp]
  \centering
  \figorplaceholder{figures/fig_real_samples.png}{8.5cm}
  \caption{Samples from the real-imagery sets with ground-truth boxes: own drone
  footage at 10, 30, and 50~m altitude, a campus scene with partial occlusion, a
  disaster-training frame from news footage of a flood, and a frame from the
  held-out unseen disaster event. Person size and scene clutter vary widely
  across the sets, which is what the adaptation must cope with.}
  \label{fig:realsamples}
\end{figure}

Zero-shot transfer is poor. On the drone test set the C2A-trained assignment
model reaches $\mathrm{AP}_{50}$ $0.335$; on the disaster R-set it reaches
$0.133$. Supervised adaptation closes most of the drone gap: fine-tuning on C2A
plus the oversampled real training sets (Section~\ref{sec:adaptmethod}) lifts the drone test to
$0.783$ with the drone data alone and to $0.795$ when the disaster and campus
sets are added, with operational recall at the fixed $0.25$ threshold rising
from $0.337$ to $0.837$. The cost on the source domain is negligible: the C2A
test score moves from $0.844$ to $0.837$ ($-0.7$ points), very-tiny recall holds
within half a point, and calibration improves (expected calibration error
$0.014 \to 0.011$). One model therefore serves the synthetic benchmark and the
real footage at the same time; the architecture, parameter count, and latency
are unchanged from the CBAM~+~P2 model measured in Table~\ref{tab:main}.
Figure~\ref{fig:sim2real} and Table~\ref{tab:multidomain} summarise the final
model across every held-out set.

\begin{figure}[tbp]
  \centering
  \figorplaceholder{figures/fig_sim2real_bars.png}{6.5cm}
  \caption{$\mathrm{AP}_{50}$ of the three model stages on the three real
  evaluation domains. Drone fine-tuning transfers to the drone test but not to
  disaster footage; adding the small disaster and campus training sets triples
  the within-video R-set score while the drone score is retained. On the unseen
  event the difference between the base and adapted models is not statistically
  significant ($p=0.47$); the middle bar is not measured there.}
  \label{fig:sim2real}
\end{figure}

\begin{table}[tbp]
  \centering
  \caption{The final adapted model on every held-out evaluation set. Operational
  recall is at the fixed confidence threshold $0.25$; the small real sets carry
  wide confidence intervals and are reported with their sizes.}
  \label{tab:multidomain}
  \begin{tabular}{@{}l r r r r@{}}
    \toprule
    Evaluation set & Images & $\mathrm{AP}_{50}$ & $F_2$ & Recall@$0.25$ \\
    \midrule
    C2A scene-disjoint test   & 2{,}040 & 0.837 & 0.824 & 0.821 \\
    Drone test                & 60      & 0.795 & 0.771 & 0.837 \\
    Disaster R-set            & 25      & 0.411 & 0.472 & 0.451 \\
    Unseen disaster event     & 23      & 0.413 & 0.451 & 0.374 \\
    Campus (occlusion)        & 25      & 0.689 & 0.666 & 0.718 \\
    \bottomrule
  \end{tabular}
\end{table}

\subsection{Disaster Case Study: Within-Video Against Cross-Event}
\label{sec:xevent}
The disaster numbers need careful reading, and the honest reading is the main
finding of this subsection. The 60 disaster training frames and the 25 R-set
evaluation frames come from the same two news videos; a perceptual-hash check
guarantees no frame is shared, but the scenes, lighting, and people overlap. On
that within-video set, adaptation works dramatically: $\mathrm{AP}_{50}$ rises
from $0.133$ to $0.411$, recall at the operating threshold from $0.127$ to
$0.451$, and recall in the 16--32~px band from $0.230$ to $0.569$. Drone
fine-tuning alone does nothing for the same set ($0.133 \to 0.124$), so the gain
comes specifically from disaster labels, not from real-image exposure in
general.

The cross-event test asks the harder question: does any of this transfer to a
disaster the model has never seen? On the 23 held-out frames from a third event,
the adapted model scores $\mathrm{AP}_{50}$ $0.413$ against the un-adapted
model's $0.396$, a gain of $+1.7$ points whose bootstrap interval spans zero
($[-3.2, 5.7]$, $p=0.47$), while precision at the operating threshold drops by
13 points ($p=0.018$). Two conclusions follow. First, fine-tuning on 60 frames
from two events adapts the model to those events rather than to the disaster
domain: the within-video gain is real but does not generalise. Second, the
synthetic-trained base model itself transfers better than expected to unseen
real disaster imagery of small and medium size ($0.396$ from composited training
data alone), and its genuine failure is concentrated below 16~px: in the R-set,
recall of sub-8-pixel people is zero for every model tested (54 instances).
Detecting the smallest people in real disaster scenes remains open, and the
within-video result bounds what a labelling budget of an afternoon can and
cannot buy.

\subsection{Statistical Validation}
\label{sec:stats}
Because several of the reported differences are small and the real test sets are
small, the key comparisons carry uncertainty estimates. A paired image-level
bootstrap resamples the test images with replacement (1{,}000 resamples), applies
the same resample to both models of a pair, and recomputes each metric, giving a
95\% percentile interval and a two-sided $p$-value for the delta; per-image
predictions are cached, so the procedure re-uses the recorded detections rather
than re-running inference. This treats variability from the finite test sample.
It does not treat training-run variability: every training run in this thesis
uses seed 0, a stated limitation of the earlier chapters. For the central
assignment claim this is being addressed directly; replicate runs of the matched
pair at two further seeds were in progress at the time of writing, and the
single-seed conclusion rests meanwhile on the tight bootstrap interval of
Table~\ref{tab:d2pair} and on the monotone dose--response of
Table~\ref{tab:alpha}, two independent lines of evidence that the very-tiny
gain is not noise.
%% TODO(seeds): when runs s3_{assign,control}_full_s{1,2} finish on PC-1, replace the
%% previous sentence with the 3-seed mean±std for VT-recall and AP_small (see
%% 27-07-2026-Novelty-Lap-4/D2_SIGNIFICANCE.md results table) and update Ch.VII limitations.

\section{Objective Achieved}
The objectives set in Chapter~I are met as follows. The additive ablation isolates
the contribution of each component and identifies CBAM~+~P2 as the configuration
that best serves the task, since it gives the best very-tiny recall and operational
$F_2$ at a parameter count below the baseline with P2 alone. The comparison in
Table~\ref{tab:sota} shows the model is competitive with published detectors on C2A
at a fraction of their size, which meets the objective of a compact,
deployment-ready detector. The state-space neck was evaluated honestly and shown not
to help, which answers the design question it was introduced to test. The extended
objectives are also met. The split audit quantified the benchmark's scene leakage
and re-established the results under a scene-disjoint protocol
(Table~\ref{tab:leak}). The label-assignment change delivered the targeted
very-tiny gain with a confidence interval that excludes zero
(Table~\ref{tab:d2pair}). The real-footage objective was met on both of its
sides: adaptation with a small labelled budget closed the drone-domain gap and
tripled the within-video disaster score, and the cross-event test measured,
rather than assumed, the limit of that adaptation
(Section~\ref{sec:xevent}).

\section{Financial Analyses and Budget}
The work used hardware already available and open-source software, so the direct
monetary cost is limited to electricity. Training the four configurations took
about 7--30 hours each on the 4070 Ti SUPER (the state-space model being the
longest at $29.5$~h), at an average board power of roughly $180$~W. The full set of
training runs therefore consumed on the order of $10$--$15$~kWh. The extended
study added roughly 120 further GPU-hours across the re-baseline, the pilot and
full-protocol assignment runs, the failed branch, the adaptation fine-tunes, and
the seed replicates, an estimated $20$~kWh at the same board power. Annotation of
the real sets used a free tier of a web labelling tool with segmentation-assisted
box drawing. No paid datasets, cloud compute, or licences were required.

\section{Conclusion}
The additive ablation shows that the gain on C2A comes almost entirely from the P2
detection scale, which lifts very-tiny recall and $\mathrm{AP}_{50}$ at a small
cost, while CBAM contributes efficiency and the state-space neck contributes
nothing at a large latency cost. The recommended CBAM~+~P2 model is competitive with
much larger published detectors on the same data. The extended study sharpened
this picture in three ways. The benchmark's scene leakage costs $1.6$ points of
$\mathrm{AP}_{50}$ and does not change any conclusion. On the smallest targets
the productive lever turned out to be label assignment rather than architecture:
the blended-similarity assigner added $2.1$ points of sub-8-pixel recall with a
confidence interval that excludes zero, while a frequency-gated branch joined
the state-space neck as a documented negative. On real footage, a labelled
budget of a few hundred frames adapted the model to its target domains at
almost no cost to the source benchmark, and a held-out unseen event showed
where that adaptation stops. The next chapters consider the wider engineering
context of the work and its conclusions.
